\section{Problem Statement}
\label{sec:problem}

\subsection{Problem}

Network proximity is a proven way to shortlist drugs for repurposing, and it has been applied to Alzheimer's disease. But the screen is expensive to run properly. It needs shortest-path distances between all pairs of $\sim$15{,}000--18{,}000 proteins, and each of the $\sim$1{,}400 or more drugs must be compared against 1{,}000 or more random samples to get a fair $z$-score. Typical implementations are sequential scripts, so researchers settle for few random samples, one version of the interaction network, and one definition of the AD module---even though results can change with each of these choices.

A second, smaller problem sits on top: choosing a good \emph{combination} of drugs from the shortlist is a combinatorial selection problem. It can be written as a QUBO, the input format for quantum optimization, but it is not yet known how well a gate-based quantum algorithm (QAOA, simulated on a GPU) handles this specific AD problem compared with the exact answer and with classical solvers.

\subsection{Proposed solution}

We propose a CUDA program that runs the full proximity screen on the GPU, built as five modules with clear inputs and outputs so each can be developed and tested separately (Table~\ref{tab:modules}).

\begin{table}[h]
\centering
\small
\begin{tabularx}{\linewidth}{@{}lXXl@{}}
\toprule
Module & Input & Output & Runs on \\
\midrule
1. Data preparation & Interaction list, drug--target list, AD gene list & Graph in compressed form (CSR), target and disease sets as node IDs & CPU \\
2. All-pairs distances & Graph (CSR) & Distance matrix $D$ ($N\times N$, 1 byte per entry) & GPU \\
3. Proximity scoring & $D$, drug target sets, AD module & Proximity and $z$-score per drug & GPU \\
4. Ranking and validation & $z$-scores, list of AD trial drugs & Ranked candidates, validation metrics & CPU \\
5. Combinations (optional) & Top-$n$ drugs, $D$ & Drug--drug separation matrix, QUBO, best combination & GPU (+ cuQuantum) \\
\bottomrule
\end{tabularx}
\caption{Pipeline modules. Module 5 is optional and independent of the graded
core (Modules 1--4).}
\label{tab:modules}
\end{table}

\subsection{Research questions}

\begin{enumerate}[label=\textbf{RQ\arabic*.},leftmargin=*]
  \item How much faster is the GPU implementation than a sequential C++ implementation for the all-pairs distance step and the proximity scoring step, and how does the speedup change with network size and number of random samples?
  \item Does the GPU implementation give exactly the same distances, and the same $z$-scores when given the same random samples, as the sequential version?
  \item Does the proximity ranking place drugs that have entered AD clinical trials higher than simple baselines (random ranking, or ranking by number of targets)?
  \item (Optional) On AD drug-combination QUBOs small enough to solve exactly, how close do classical heuristics and simulated QAOA get to the true best answer under the same budget, and how do biology-derived instances compare with random ones? (Appendix~\ref{app:extension})
\end{enumerate}

\subsection{Objectives}

\paragraph{General objective.} To build and evaluate a GPU-accelerated
network proximity screen for Alzheimer's disease drug repurposing.

\paragraph{Specific objectives.}
\begin{enumerate}[nosep]
  \item Implement all-pairs shortest-path distances on the GPU using bit-parallel multi-source BFS.
  \item Implement proximity and $z$-score computation for all drugs on the GPU, with random sampling done in parallel.
  \item Implement the same steps sequentially in C++ and compare execution times.
  \item Verify that GPU and CPU results match.
  \item Rank AD repurposing candidates and measure how well the ranking recovers drugs already in AD trials.
  \item (Optional) Formulate AD drug-combination selection as a QUBO and compare GPU brute force, simulated annealing, and QAOA.
\end{enumerate}

\subsection{Scope and limitations}

\begin{itemize}[nosep]
  \item Public data only: interaction, drug--target, and AD gene databases. No patient data are used.
  \item Single GPU. The network has up to about 20{,}000 nodes, so the full distance matrix (at most 400~MB at one byte per entry) fits in GPU memory.
  \item Results are \emph{hypotheses}, not treatment recommendations. Closeness in a network does not prove a drug works; clinical evidence is outside the project's scope.
  \item The interactome is incomplete and biased toward well-studied proteins; we reduce this by comparing more than one database.
  \item The quantum part is simulation only, limited to about 20--30 drugs (one qubit each). It does not claim to be faster than classical methods; its purpose is to measure solution quality.
\end{itemize}

\subsection{Significance}

A fast screen lets researchers run the checks that are currently skipped: more random samples for stable $z$-scores, several networks, and several definitions of the AD module (for example, early-stage vs.\ late-stage genes). For the course, it is a clear case study of turning a graph problem into bitwise, data-parallel GPU work. For the optional part, it provides a realistic, biology-derived benchmark for quantum optimization.
